% Generated draft: requires mathematical and source review.
\documentclass[11pt]{article}
\usepackage[margin=25mm]{geometry}
\usepackage{fontspec}
\setmainfont{Latin Modern Roman}
\usepackage{amsmath,amssymb,mathtools,unicode-math}
\setmathfont{Latin Modern Math}
\usepackage{enumitem,needspace,xurl,hyperref}
\setlength{\parindent}{0pt}
\setlength{\parskip}{5pt}
\setlength{\emergencystretch}{3em}
\newcommand{\N}{\mathbb{N}}
\newcommand{\Z}{\mathbb{Z}}
\newcommand{\Q}{\mathbb{Q}}
\newcommand{\R}{\mathbb{R}}
\newcommand{\C}{\mathbb{C}}
\newcommand{\F}{\mathbb{F}}
\newcommand{\A}{\mathbb{A}}
\newcommand{\PP}{\mathbb{P}}
\newcommand{\E}{\mathbb{E}}
\newcommand{\eps}{\varepsilon}
\newcommand{\dd}{\,\mathrm{d}}
\begin{document}
\section*{UnsolvedMath ID: 30001658}
\textbf{Source problem code:} catalogue entry
\subsection{Definitions and mathematical statement}
Let \(Q_n=\{0,1\}^n\) be the Boolean cube. Its dimension must be a nonnegative integer; the supplied record does not explicitly specify the range of \(n\), including whether dimension zero is intended. For \(x,y\in Q_n\), interpret \(x\sim y\) as meaning that \(x\) and \(y\) differ in exactly one coordinate. Let
\[
\mathcal E_n=\bigl\{\{x,y\}:x,y\in Q_n,\ x\sim y\bigr\}.
\]
For \(f:Q_n\to\R\), the two possible edge-counting conventions are
\[
E_{\mathrm{unord}}(f,f)=\sum_{\{x,y\}\in\mathcal E_n}(f(x)-f(y))^2,
\qquad
E_{\mathrm{ord}}(f,f)=\sum_{x\in Q_n}\sum_{\substack{y\in Q_n\\y\sim x}}(f(x)-f(y))^2
=2E_{\mathrm{unord}}(f,f).
\]
The supplied record writes \(E(f,f)=\sum_{x\sim y}(f(x)-f(y))^2\) without specifying which convention is intended. This ambiguity affects the proposed inequality and remains unresolved.

\textbf{Conjecture as supplied.} For every intended dimension \(n\), every nonempty subset \(A\subseteq Q_n\), and every function \(f:Q_n\to\R\),
\[
E(f,f)+4\sum_{x\in Q_n}f(x)^2
\geq
2\left(
\frac{\log_2(2^n/|A|)}{|A|}
\sum_{x\in A}|f(x)|
\right)^2.
\]
Here \(|A|\) is the cardinality of \(A\), \(\log_2\) is the base-two logarithm, and the displayed vertex sums are unnormalized sums. The supplied formulation imposes no support condition on \(f\), no normalization of \(f\), and no further restriction on the nonempty set \(A\). No such hypothesis has been added. A definitive formulation requires resolution of the edge-counting convention and confirmation of the intended dimension range.

The registry labels this a conjecture with status \texttt{open}. Its supplied literature assessment says that known logarithmic-Sobolev and Talagrand inequalities establish only weaker or supported cases, without specifying those cases. These are supplied source qualifications, not independently verified mathematical or literature conclusions.
\subsection{Short English statement}
The supplied conjecture gives a lower bound for a Boolean-cube energy plus four times the squared \(\ell^2\) norm, using the average of \(|f|\) on an arbitrary nonempty subset and its base-two logarithmic size deficit. The edge-counting convention and the dimension range require clarification. This is a statement-only draft requiring human mathematical review.
\subsection{Sources}
\begin{itemize}
\item Supplied local UnsolvedMath registry record, permanent ID \(30001658\), title ``Functional Inequality on the Boolean Cube,'' exact source problem code \texttt{catalogue entry}, in the collection ``Oberwolfach Reports Open Problems.'' The supplied proposed year is 2011; no individual proposer is identified.
\item Supplied report citation: \emph{Combinatorics} (2011), Oberwolfach Reports, pp.\ 5--83, DOI \texttt{10.4171/owr/2011/01}. Literature-source label: ``Oberwolfach Report 01/2011.'' \url{https://doi.org/10.4171/owr/2011/01}. The background describes these reports as documenting mathematical workshops and their open-problem discussions; no precise page or passage for this conjecture is supplied.
\item Supplied literature reference labeled ``Samorodnitsky'': \url{https://arxiv.org/abs/0807.1679}. The record supplies no title, theorem number, or precise explanation of its applicability.
\item Literature checked at: \texttt{2026-08-21}. The supplied dated curation labels the problem open and describes the sharp functional inequality as a source conjecture, with known logarithmic-Sobolev and Talagrand inequalities providing only weaker or supported cases. The background explicitly requires independent verification of time-sensitive status and cited claims. This assessment is supplied dated curation and has not been independently verified here.
\item The supplied statement status is \texttt{corrected\_verified}, accompanied by the assertion that the official report was checked and the statement repaired to match its source passage. This is a registry verification claim, not verification performed for this draft.
\end{itemize}
\end{document}
